%% ma: L11-B02-0003
%% lop: 11
%% chuong: 1
%% bai: 2
%% dang: 11.2.3
%% loai: TLN
%% muc: VD
%% dap_an: "-0,9"
%% nguon: "de-cuong-giua-ky-1-toan-11-nam-2026-2027-truong-thpt-hoang-van-thu-ha-noi-2-8.pdf (M08, Đề cương giữa kì 1 Toán 11 THPT Hoàng Văn Thụ, Hà Nội), Phần 3 Câu 2"
%% kien_thuc: "Bài 1 (giá trị lượng giác của góc lượng giác), Bài 2 (công thức lượng giác)"
%% trang_thai: cho-duyet
%% ly_do: "Dạng mới đề xuất 11.2.3; Có sửa đề so với gốc (xem ghi_chu)"
%% ghi_chu: "Gốc ghi gọn \"π/2 < α, β < π\" và không yêu cầu làm tròn; đáp số chính xác −(2 + 2√10)/9 ≈ −0,925 không ghi được trong 4 ký tự. Bản kho bổ sung \"kết quả làm tròn đến hàng phần mười\"."
\begin{ex}
Cho các góc $\alpha$, $\beta$ thỏa mãn $\dfrac{\pi}{2}<\alpha<\pi$, $\dfrac{\pi}{2}<\beta<\pi$, $\sin\alpha=\dfrac13$, $\cos\beta=-\dfrac23$. Tính $\sin(\alpha+\beta)$ (kết quả làm tròn đến hàng phần mười).
\shortans{-0,9}
\loigiai{$\cos\alpha=-\sqrt{1-\dfrac19}=-\dfrac{2\sqrt2}{3}$, $\sin\beta=\sqrt{1-\dfrac49}=\dfrac{\sqrt5}{3}$.

$\sin(\alpha+\beta)=\dfrac13\cdot\left(-\dfrac23\right)-\dfrac{2\sqrt2}{3}\cdot\dfrac{\sqrt5}{3}=-\dfrac{2+2\sqrt{10}}{9}\approx-0{,}9$.}
\end{ex}
